\subsection{Witt's Theorem.}

Given two vector spaces E, E' over A equipped respectively with two sesquilinear forms \(\Phi\) , \(\Phi'\) (resp. with two quadratic forms Q, Q'), one calls a metric homomorphism from E to E' any linear map \(u\) from E to E' such that \(\Phi'(u(x), u(y)) = \Phi(x, y)\) (resp. Q'(u(x)) = Q(x)) for \(x \in E\) , \(y \in E\) . If E and E' have the same finite dimension and if \(\Phi\) (resp. Q) is nondegenerate, every metric homomorphism \(u\) from E to E' is an isomorphism, since \(u(x) = 0\) implies \(\Phi(x, y) = 0\) for all \(y \in E\) , hence \(x = 0\) ; thus \(u\) is injective, hence bijective since E and E' have the same finite dimension.

THEOREM 1 (Witt). — Let E and E' be two finite-dimensional vector spaces, equipped respectively with two nondegenerate ε-hermitian forms Φ and Φ' satisfying condition (T) of no. 2 (resp. with two nondegenerate quadratic forms Q and Q'), and isomorphic for these structures. Given any subspace F of E, every injective metric homomorphism of F into E' extends to a metric isomorphism of E onto E'.

Using the given isomorphism from E onto \(\mathbf{E}'\) , one sees that it suffices to show that every injective metric homomorphism \(u\) from F into E extends to a metric automorphism of E. Note that if, for \(i = 1, 2\) , \(\mathbf{F}_i\) is a subspace of E and \(u_i\) a metric homomorphism from \(\mathbf{F}_i\) into E, such that \(\mathbf{F}_1 \cap \mathbf{F}_2 = \{0\}\) and \(\Phi(u_1(x_1), u_2(x_2)) = \Phi(x_1, x_2)\) for \(x_i \in \mathbf{F}_i\) ( \(i = 1, 2\) ), then the homomorphism \(\nu : x_1 + x_2 \to u_1(x_1) + u_2(x_2)\) of \(\mathbf{F}_1 + \mathbf{F}_2\) into E which extends \(u_1\) and \(u_2\) is metric: indeed, for any \(x_i, y_i\) in \(\mathbf{F}_i\) ( \(i = 1, 2\) ), the expansion of each of the expressions \(\Phi(x_1 + x_2, y_1 + y_2)\) and \(\Phi(u_1(x_1) + u_2(x_2), u_1(y_1) + u_2(y_2))\) (resp. \(\mathrm{Q}(x_1 + x_2)\) and \(\mathrm{Q}(u_1(x_1) + u_2(x_2))\) ) contains four (resp. three) terms equal each to each according to the hypotheses made. Moreover, if \(u_1\) and \(u_2\) are injective and if \(u_1(\mathbf{F}_1) \cap u_2(\mathbf{F}_2) = \{0\}\) , then \(\nu\) is injective.

\begin{enumerate}
\def\labelenumi{\arabic{enumi})}
\tightlist
\item
  Let us first prove Witt's theorem in the case where the set of fixed points of u is a hyperplane U of F. The set of vectors of the form \(u(x) - x\) with \(x \in F\) is then a line D. If \(F'\) is a subspace orthogonal to D such that \(F' \cap F = F' \cap u(F) = \{0\}\) , one will have \(\Phi(u(x), y) = \Phi(x, y)\) for \(x \in F\) and \(y \in F'\) ; our initial remark therefore applies to u and to the identity map of \(F'\) into E, showing that u extends to \(F + F'\) while leaving the points of \(F'\) fixed; the set of vectors of the form \(u(x) - x (x \in F + F')\) is still the line D. Now we have, for \(x \in F, y \in F\)
\end{enumerate}

\[
\Phi (u (x), u (y) - y) = \Phi (u (x), u (y)) - \Phi (u (x), y) = \Phi (x - u (x), y),
\]

which, when \(x \in \mathrm{U}\) (that is, when \(u(x) = x\) ), shows that \(x \in \mathrm{D}^0\) ; in other words, we have \(\mathrm{U} \subset \mathrm{D}^0\) . We distinguish two cases:

\begin{enumerate}
\def\labelenumi{\alph{enumi})}
\item
  \(F \subset D^{0}\) . Formula (5) shows that \(u(F)\) is not contained in \(D^{0}\) , hence \(F \cap D^{0} = u(F) \cap D^{0} = U\) . One can then take for \(F'\) a supplement of U in \(D^{0}\) ; since \(F + F'\) contains the hyperplane \(D^{0}\) and is distinct from it, we have \(F + F' = E\) , and we have found in this case the desired extension of u to E.
\item
  \(\mathbf{F} \subset \mathbf{D}^0\) . Formula (5) shows that \(u(\mathbf{F}) \subset \mathbf{D}^0\) , and hence that
\end{enumerate}

D ⊂ D⁰ ; the line D is therefore isotropic (resp. singular, since we have Q(u(x) - x) = Q(u(x)) - Φ(x, u(x)) + Q(x) = 2Q(x) - Φ(x, x) = 0 for x ∈ F). We shall show that, under these conditions, there exists a subspace F' of D⁰ which is a supplement of F and of u(F) in D⁰. This is immediate if F = u(F). Otherwise, let x and y be vectors such that x ∈ F, x ∉ U, y ∈ u(F), y ∉ U; we then have F = U + Ax, u(F) = U + Ay, and F does not contain x + y, otherwise y = (x + y) - x would belong to F ∩ u(F) = U; we see similarly that x + y does not belong to u(F); thus the line A(x + y) is a supplement of F and of u(F) in the subspace F + u(F); it then suffices to set F' = A(x + y) + G, where G is a supplement of F + u(F) in D⁰. This being so, we have F + F' = u(F) + F' = D⁰, and, in this case, what was said at the beginning of 1) shows that there exists an extension of u to the hyperplane D⁰ of E, and that D⁰ is stable for this extension.

We are therefore reduced to the case where F is the hyperplane D⁰ and where u is an automorphism of F. Let us prove that, for every z ∈ E, there exists z' ∈ E such that

\[
\Phi (u (x), z ^ {\prime}) = \Phi (x, z)\tag{6}
\]

for all \(x \in \mathbf{F}\) ; indeed the linear form \(x \to \Phi(u^{-1}(x), z)\) on \(\mathbf{F}\) is the restriction of a linear form on \(\mathbf{E}\) , a form which is of the type \(x \to \Phi(x, z')\) since \(\Phi\) is nondegenerate; hence (6) is valid. Moreover, if \(z \notin \mathbf{F}\) there exists a vector \(z' \in \mathbf{E}\) satisfying (6) and such that \(\Phi(z', z') = \Phi(z, z)\) (resp. \(Q(z') = Q(z)\) ). Indeed, formula (6) remains valid if one adds to \(z'\) an element \(u(y) - y\) ( \(y \in \mathbf{F}\) ) of \(\mathbf{D}\) since \(\mathbf{F} = \mathbf{D}^0\) , and Lemma 1 of No. 2 allows us to conclude since \(z\) is not orthogonal to \(\mathbf{D}\) Our initial remark then shows that there exists a metric homomorphism \(v\) of \(\mathbf{F} + A z = E\) into \(E\) which extends \(u\) and which transforms \(z\) to \(z'\) . Since \(\Phi\) is nondegenerate, \(v\) is the metric automorphism of \(E\) sought.

\begin{enumerate}
\def\labelenumi{\arabic{enumi})}
\setcounter{enumi}{1}
\tightlist
\item
  In the general case, we will argue by induction on \(r = \dim F\) The case \(r = 0\) is trivial. Let then \(r > 0\) , that is to say \(F \neq \{0\}\) and let U be a hyperplane of F. The restriction \(u_0\) of \(u\) to U extends, by the induction hypothesis, to a metric automorphism \(v_0\) of E. If \(v_0\) extends \(u\) , the theorem is proved. Otherwise, U is the set of invariant elements by \(v_{0}^{-1}u\) , and there exists, by 1), a metric automorphism \(v_{1}\) of E extending \(v_{0}^{-1}u\) . The automorphism \(v_{0}v_{1}\) is then the desired extension of u. QED.
\end{enumerate}

COROLLARY 1. --- Let, for \(i=1,2\) , \(\mathbf{E}_{i}\) be a finite-dimensional vector space, \(\Phi_{i}\) a nondegenerate \(\varepsilon\) -hermitian form on \(\mathbf{E}_{i}\) satisfying (T) (resp. \(\mathbf{Q}_{i}\) a nondegenerate quadratic form on \(\mathbf{E}_{i}\) ), \(\mathbf{E}_{i}^{\prime}\) and \(\mathbf{E}_{i}^{\prime\prime}\) two orthogonal subspaces of \(\mathbf{E}_{i}\) of which \(\mathbf{E}_{i}\) is the direct sum. If the forms \(\Phi_{1}\) and \(\Phi_{2}\) (resp. \(\mathbf{Q}_{1}\) and \(\mathbf{Q}_{2}\) ) are equivalent, and if their restrictions to \(\mathbf{E}_{1}^{\prime}\) and \(\mathbf{E}_{2}^{\prime}\) are equivalent, then so are their restrictions to \(\mathbf{E}_{1}^{\prime\prime}\) and \(\mathbf{E}_{2}^{\prime\prime}\) .

Indeed, let \(u\) be a metric isomorphism of \(\mathbf{E}_{1}^{\prime}\) onto \(\mathbf{E}_{2}^{\prime}\) . By th. 1, \(u\) extends to a metric isomorphism \(\nu\) of \(\mathbf{E}_{1}\) onto \(\mathbf{E}_{2}\) . Since \(\Phi_{i}\) is nondegenerate, \(\mathbf{E}_{i}^{\prime\prime}\) is the orthogonal of \(\mathbf{E}_{i}^{\prime}\) in \(\mathbf{E}_{i}\) , hence \(\nu\) maps \(\mathbf{E}_{1}^{\prime\prime}\) on \(\mathbf{E}_{2}^{\prime\prime}\) . QED.

COROLLARY 2. --- Under the hypotheses of Theorem 1, the group of metric automorphisms of E acts transitively on the totally isotropic (resp. totally singular) subspaces of E of a given dimension. Moreover, if F is a totally isotropic (resp. totally singular) subspace of E, then every bijective linear map from F onto F is induced by a metric automorphism of E.

COROLLARY 3. --- Let Q be a nondegenerate quadratic form on a finite-dimensional vector space E over an algebraically closed field A. The group of metric automorphisms of E acts transitively on the non-isotropic subspaces of E of a given dimension.

This follows immediately from th. 1 and cor. 2 of prop. 3.

Exercises. -- 1) a) Let K be a field of characteristic 2, J : \(\xi \to \bar{\xi}\) an involutive antiautomorphism of K, Z the center of K. Show that if the restriction of J to Z is not the identity, every element \(\mu\) of K such that \(\bar{\mu} = \mu\) is of the form \(\lambda + \bar{\lambda}\) (note that there is an element \(\rho \neq 0\) in Z which can be written in the form \(\zeta + \bar{\zeta}\) , with \(\zeta \in \mathbb{Z}\) ); every Hermitian form on a vector space over K then satisfies condition (T).

\begin{enumerate}
\def\labelenumi{\alph{enumi})}
\setcounter{enumi}{1}
\tightlist
\item
  Give examples of fields of characteristic 2 admitting an involutive antiautomorphism \(\xi \to \overline{\xi}\) distinct from the identity map, and for which there exist elements \(\mu = \bar{\mu}\) which are not of the form \(\lambda + \bar{\lambda}\) (cf. Chap. VIII, § 11, Exerc. 4).
\end{enumerate}

\begin{enumerate}
\def\labelenumi{\arabic{enumi})}
\setcounter{enumi}{1}
\tightlist
\item
  Let A be a field, E a vector space over A, \(\Phi\) (resp. Q) a nondegenerate \(\varepsilon\) -hermitian form on E, satisfying condition (T) (resp. a nondegenerate quadratic form on E, \(\Phi\) then denoting the symmetric bilinear form associated with Q).
\end{enumerate}

\begin{enumerate}
\def\labelenumi{\alph{enumi})}
\item
  Show that for a plane P ⊂ E to be isotropic (resp. singular) and not totally isotropic (resp. not totally singular), it is necessary and sufficient that it contain only one isotropic (resp. singular) line (cf. exerc. 14 e)).
\item
  We assume that dim \(E \geqslant 3\) and that there exist isotropic vectors in E \(\neq 0\) Show that if P is a non-totally isotropic plane in E, there exists a non-isotropic subspace \(V \subset E\) of dimension 3, containing isotropic vectors \(\neq 0\) , and such that \(P \subset V\) .
\end{enumerate}

\begin{enumerate}
\def\labelenumi{\arabic{enumi})}
\setcounter{enumi}{2}
\item
  With the same hypotheses as in exerc. 2, show that if dim E ≥ 3, every isotropic line in E is the intersection of two non-isotropic planes.
\item
  The hypotheses are those of exerc. 2, and it is further assumed that E is of finite dimension.
\end{enumerate}

\begin{enumerate}
\def\labelenumi{\alph{enumi})}
\item
  If the index v of \(\Phi\) (resp. Q) is \(\geqslant 1\) , show that for every isotropic (resp. singular) vector \(a \neq 0\) in E, there exists a basis \((e_i)\) of E consisting of isotropic (resp. singular) vectors, such that \(e_1 = a\) (cf. exerc. 14 e)).
\item
  Let V, W be two totally isotropic (resp. totally singular) subspaces of the same dimension \(r \leqslant v\) ; show that there exist two maximal totally isotropic (resp. totally singular) subspaces \(V_{1}\) , \(W_{1}\) , such that \(V \subset V_{1}\) , \(W \subset W_{1}\) , and \(V_{1} \cap W_{1} = V \cap W\) . (If U = V ∩ W, reason in \(U^{0}/U\) ).
\item
  Let V, W, \(V_{1}\) , \(W_{1}\) be four totally isotropic (resp. totally singular) subspaces of the same dimension, such that \(V + W\) and \(V_{1} + W_{1}\) are non-isotropic. Show that there exists a metric automorphism u of E such that \(u(V) = V_{1}\) and \(u(W) = W_{1}\) .
\item
  Let \(f\) be a linear form on E, \(\alpha\) an element of A of the form \(\lambda + \varepsilon\bar{\lambda}\) (resp. an element of A). We consider the sesquilinear form on E
\end{enumerate}

\[
(x, y) \rightarrow \Phi_ {1} (x, y) = \Phi (x, y) + f (x) \alpha \overline {{f (y)}}
\]

(resp. the quadratic form

\[
x \rightarrow \mathrm{Q} _ {1} (x) = \mathrm{Q} (x) + \alpha (f (x)) ^ {2}).
\]

Show that if \(\Phi_{1}\) (resp. \(Q_{1}\) ) is nondegenerate and if \(\nu_{1}\) denotes its index, we have \(|\nu_{1} - \nu| \leqslant 1\) .

\begin{enumerate}
\def\labelenumi{\arabic{enumi})}
\setcounter{enumi}{4}
\tightlist
\item
  \begin{enumerate}
  \def\labelenumii{\alph{enumii})}
  \tightlist
  \item
    Let B be a ring, \(\xi \to \overline{\xi}\) an involutive antiautomorphism of B, \(\varepsilon\) an element of the center of B such that \(\varepsilon \varepsilon = 1\) . Show that if \(\beta\) is an invertible element of B such that \(\beta + \varepsilon \overline{\beta} \neq 0\) , there exists an invertible element \(\mu \neq 1\) in B, such that \(\mu (\beta + \varepsilon \overline{\beta}) \mu = \beta + \varepsilon \overline{\beta}\) (Show that one can take \(\mu\) such that \(\mu \beta \overline{\mu} = \beta\) ).
  \end{enumerate}
\end{enumerate}

\begin{enumerate}
\def\labelenumi{\alph{enumi})}
\setcounter{enumi}{1}
\tightlist
\item
  Let A be a field, and let E be a vector space over A, \(\Phi\) a nondegenerate ε-hermitian sesquilinear form on E, satisfying (T). Show that if Φ is not alternating, then for every non-isotropic hyperplane H of E, there exists a metric automorphism of E, distinct from the identity, leaving every element of H invariant (use a)).
\end{enumerate}

\begin{enumerate}
\def\labelenumi{\arabic{enumi})}
\setcounter{enumi}{5}
\item
  Given the hypotheses of exercise 2, let \(a\) be an isotropic vector \(\neq 0\) in E (resp. a nonsingular isotropic vector (note that such vectors exist only if A has characteristic 2)). Let \(\lambda \in \mathbf{A}\) such that \(\lambda + \varepsilon \bar{\lambda} = 0\) (resp. \(\lambda = (\mathrm{Q}(a))^{-1}\) ); show that the transvection \(x \to x + \Phi(x, a) \lambda a\) (chap. II, § 6, exerc. 7) is a metric automorphism of E; conversely.
\item
  Under the hypotheses of Exercise 2, let G be the group of metric automorphisms of E. Show that the only semilinear bijections of E onto itself that commute with every element of G are the homotheties of E, except in the following three cases: dim E = 2, G is the group of metric automorphisms corresponding to a quadratic form of index 1 on E, and A is one of the three fields \(\mathbf{F}_{2}, \mathbf{F}_{3}\) or \(\mathbf{F}_{4}\) (Use Exercises 5, 6, and 3; examine separately the case of a quadratic form on a vector space of dimension 2).
\end{enumerate}

*8) Let A be a field, E a vector space of finite dimension \textgreater0 over A, Φ a nondegenerate ε-hermitian sesquilinear form on E, satisfying (T). Let M(Φ) be the group of multipliers of the similarities of E for Φ (§ 6, no. 5).

\begin{enumerate}
\def\labelenumi{\alph{enumi})}
\item
  Let \(V_{1}\) , \(V_{2}\) two vector subspaces of E, of the same dimension, and let \(\Phi_{1}\) , \(\Phi_{2}\) be the restrictions of \(\Phi\) to \(V_{1}\) , \(V_{2}\) respectively. For there to exist a similarity u such that \(u(V_{1}) = V_{2}\) , it is necessary and sufficient that there exist \(\alpha \in M(\Phi)\) such that \(\Phi_{2}\) is equivalent to \(\alpha\Phi_{1}\) (use Witt's theorem).
\item
  Let (F, F', G) be a Witt decomposition of E (no. 2), and let \(\Phi_0\) be the restriction of \(\Phi\) to the non-isotropic subspace G. Show that one has M( \(\Phi\) ) = M( \(\Phi_0\) ) if G ≠ \{0\}. (Use Witt's theorem and prop. 2 of no. 2).
\item
  Show that if the index v of Φ is such that dim E = 2v, M(Φ) is the group of elements ζ ≠ 0 of the center of A such that ζ̄ = ζ. If dim E = 2v + 1, M(Φ) is the group of elements of the form ρρ̄, where ρ runs through the multiplicative group of elements ≠ 0 of the center of A (use Witt's theorem). (Cf. § 10, exerc. 18).*
\end{enumerate}

*9) Let A be a commutative field, E a vector space over A, Q a nondegenerate quadratic form on E. One calls a similarity for Q any automorphism \(u\) of E such that there exists an element \(\alpha \neq 0\) of A for which \(Q(u(x)) = \alpha Q(x)\) for every \(x \in E\) ; \(u\) is then also a similarity for the bilinear form associated with Q. Assuming the dimension of E finite and \(>0\) , state and prove for the similarities relative to Q the analogues of the results of exerc. 8.*

*10) Let A be a field, E a vector space over A of dimension \(\geqslant 2\) , \(\Phi_1\) (resp. \(\Phi_2\) ) a nondegenerate sesquilinear form \(\varepsilon_1\) -hermitian (resp. \(\varepsilon_2\) -hermitian) on E for an involutive antiautomorphism \(J_1\) (resp. \(J_2\) ) of A, satisfying condition (T). Show that if the group of metric automorphisms of E for \(\Phi_1\) is a subgroup of the group of similarities for \(\Phi_2\) , there exists \(\alpha \in A\) such that \(\Phi_2 = \Phi_1\alpha\) (use Exercises 5 b) and 6).

Prove the analogous property when A is assumed commutative, and that \(\Phi_{1}\) and \(\Phi_{2}\) are replaced in the statement by two nondegenerate quadratic forms \(Q_{1}, Q_{2}\) on E.*

\begin{enumerate}
\def\labelenumi{\arabic{enumi})}
\setcounter{enumi}{10}
\tightlist
\item
  Let A be a ring, J an involutive antiautomorphism of A, and E an A-module admitting a finite basis \((e_{i})\) , \(\Phi\) a nondegenerate \(\varepsilon\) -hermitian form on E, R the matrix of \(\Phi\) with respect to \((e_{i})\) ; the group of metric automorphisms of \(\Phi\) is identified with the group G of invertible matrices U such that \(^{t}U.R.U^{j}=R\) .
\end{enumerate}

\begin{enumerate}
\def\labelenumi{\alph{enumi})}
\tightlist
\item
  Assume that there exists a matrix \(P\) such that \(R = {}^t P + \varepsilon P^J\) . Show that for every matrix \(S\) such that \({}^t S + \varepsilon S^J = 0\) , and that \(P + S\) is invertible, \(U = (^{t}P^{J} - \varepsilon^{J}S)^{-1}(P + S)\) belongs to G, and \(\varepsilon I + U\) is invertible. Conversely (show that for every matrix \(U \in G\) such that \(\varepsilon I + U\) is invertible, we have
\end{enumerate}

\[
\varepsilon (\varepsilon I + ^ {t} U) ^ {- 1} R + \varepsilon^ {\mathrm{J}} R (\varepsilon^ {\mathrm{J}} I + U ^ {\mathrm{J}}) ^ {- 1} = R).
\]

\begin{enumerate}
\def\labelenumi{\alph{enumi})}
\setcounter{enumi}{1}
\tightlist
\item
  Show that the condition of a) is satisfied when \(\Phi\) satisfies condition (T). Case where in A the equation \(2\xi = \alpha\) admits one and only one solution for every \(\alpha \in \mathbf{A}\) .
\end{enumerate}

¶ 12) Let A be a commutative field, E a vector space of finite dimension n over A, Φ a nondegenerate ε-hermitian sesquilinear form on E.

\begin{enumerate}
\def\labelenumi{\alph{enumi})}
\item
  Let \(u\) be an endomorphism of E; show that if the \(r_i\) ( \(1 \leqslant i \leqslant m\) ) are the similarity invariants of \(u\) (chap. VII, § 5, no. 1, def. 1), the similarity invariants of the adjoint \(u^*\) of \(u\) with respect to \(\Phi\) are the polynomials \(\bar{r}_i\) ( \(1 \leqslant i \leqslant m\) ), where \(\bar{r}_i\) is deduced from \(r_i\) by applying to each coefficient the automorphism J (cf. chap. VII, § 5, exerc. 2). For every monic irreducible polynomial \(p \in \mathrm{A}[\mathrm{X}]\) dividing the minimal polynomial of \(u\) , let \(\mathrm{F}_k(u, p)\) be the kernel of \((p(u))^k\) in E, and let \(\mathrm{F}(u, p)\) be the union of the \(\mathrm{F}_k(u, p)\) for all integers \(k > 0\) . Show that if \(p\) and \(q\) are two distinct monic irreducible polynomials dividing the minimal polynomial of \(u\) , the subspaces \(\mathrm{F}(u, p)\) and \(\mathrm{F}(u^*, \bar{q})\) are orthogonal (use Bezout's identity). Finally, if G is a vector subspace of E such that \(u(\mathrm{G}) \subset \mathrm{G}\) , one has \(u^*(\mathrm{G}^0) \subset \mathrm{G}^0\) .
\item
  We assume that \(uu^{*} = u^{*}u\) (the case where one says that \(u\) is a normal endomorphism for \(\Phi\) ; cf.§ 7, no. 3); show that one then has \(u^{*}(F_{k}(u,p)) \subset F_{k}(u,p)\) for all \(k\) and consequently \(u^{*}(F(u,p)) \subset F(u,p)\) . If one sets \(G(p,\bar{q}) = F(u,p) \cap F(u^{*},\bar{q})\) , show that E is the direct sum of the subspaces \(G(p,\bar{q})\) , and that \(G(p,\bar{q})\) and \(G(p_{1},\bar{q}_{1})\) are orthogonal if \(p \neq q_{1}\) or if \(p_{1} \neq q\) ; in particular \(G(p,\bar{q})\) is totally isotropic if \(p \neq q\) . Show that \(G(p,\bar{p})\) is reduced to 0 or non-isotropic, and that if \(p \neq q\) , no nonzero vector of \(G(p,\bar{q})\) is orthogonal to \(G(q,\bar{p})\) (using the fact that \(\Phi\) is nondegenerate); deduce that if \(p \neq q\) , \(G(p,\bar{q})\) and \(G(q,\bar{p})\) are totally isotropic subspaces of the same dimension, and that \(G(p,\bar{q}) + G(q,\bar{p})\) is non-isotropic.
\item
  Assume that J is not the identity or that A does not have characteristic 2, and that \(u^{*} = u\) . Let M be the set of non-isotropic subspaces \(\mathbf{M} \subset \mathrm{G}(p, \overline{p})\) stable under \(u\) (hence submodules of the A{[}X{]}- module \(E_{u}\) (chap. VII, § 5, no. 1)). Show that if M is a minimal element of M, then M is an indecomposable submodule of \(E_{u}\) (chap. VII, § 4, no. 7). (Assume that M is a direct sum of an indecomposable submodule \(M_{1}\) and of a submodule \(M_{2} \neq \{0\}\) , the minimal polynomials \(p^{h}\) and \(p^{k}\) of the restrictions of u to \(M_{1}\) and \(M_{2}\) respectively being such that \(h \geqslant k\) . Note then that \(M_{1}\) is necessarily isotropic and that every \(z \neq 0\) in \(M_{1}\) such that \(p(u).z = 0\) is orthogonal to \(M_{1}\) (using the fact that every submodule of \(M_{1}\) is monogenic); writing that \(z = (p(u))^{h-1}.x\) and that z is not orthogonal to \(M_{2}\) , and deduce that one necessarily has k = h. Then show that there exists an indecomposable submodule \(N_{2}\) of \(M_{2}\) such that \(p^{h}\) be the minimal polynomial of the restriction of u to \(N_{2}\) , and that \(M_{1} + N_{2}\) be non-isotropic; conclude from this that \(M_{2} = N_{2}\) . Finally, if \(y \in M_{2}\) is not orthogonal to z, consider the submodule P of M generated by \(w = x + \lambda y\) , where \(\lambda \in A\) , and show that one can take \(\lambda\) such that P is non-isotropic, by proving that one has \(\Phi((p(u))^{h-1}.w,w) \neq 0\) ; which leads to a contradiction).
\end{enumerate}

\(d)\) Deduce from \(c)\) that \(\mathrm{G}(p,\bar{p})\) is a direct sum of indecomposable submodules \(\mathrm{H}_{i}\) pairwise orthogonal. If \(p^{h}\) is the minimal polynomial of the restriction of \(u\) to \(\mathrm{H}_{i}\) , and if \(d\) is the degree of \(p\) , show that there exists in \(\mathrm{H}_{i}\) a totally isotropic subspace of dimension \(d\) . {[}h/2{]}. Case where E contains no isotropic vector \(\neq 0\) (cf.§ 7, no. 3).

\begin{enumerate}
\def\labelenumi{\alph{enumi})}
\setcounter{enumi}{4}
\item
  State and prove the properties analogous to those of c) and d) when one has \(u^{*} = -u\) or \(u^{*}u = 1\) .
\item
  Give an example where n = 4, Φ is symmetric and of index 2, p = \(\bar{p} = X - 1\) , u is normal, E = G(p, \(\bar{p}\) ), but E is not a direct sum of minimal submodules of M, and where there exists an eigenvector of u that is not an eigenvector of \(u^{*}\) (cf.§ 7, no. 3).
\end{enumerate}

\begin{enumerate}
\def\labelenumi{\arabic{enumi})}
\setcounter{enumi}{12}
\item
  The hypotheses are those of Exercise 2, and one further assumes that E admits a countable basis \((e_n)\) . Let F be a totally isotropic (resp. totally singular) subspace of E such that \(F^{00} = F\) . Show that there exists a totally isotropic (resp. totally singular) subspace \(F'\) such that: \(1^\circ F \cap F' = \{0\}\) ; \(2^\circ\) there exists a basis \((a_m)_{m \in I}\) of F and a basis \((a'_m)_{m \in I}\) of \(F'\) (I interval of N starting at 0) such that \(\Phi(a_i, a_j') = \delta_{ij}\) for every pair of indices; \(3^\circ (F + F')^{00} = F + F'\) and E is the direct sum of \(F + F'\) and of \(G = (F + F')^0\) . (Construct by induction an increasing sequence \((L_n)\) of non-isotropic subspaces, with union E, such that dim \(L_{n+1} - \dim L_n = 2\) , and apply prop. 2 of no. 2 to each of the \(L_n\) ; to construct this sequence, one will consider, for every \(n\) , the smallest integer \(k\) such that \(e_k \notin L_n\) , and one will use exerc. 9 b) of § 1).
\item
  Let A be a field of characteristic 2, E a vector space of finite dimension n over A, Φ a nondegenerate Hermitian form on E, not necessarily satisfying condition (T).
\end{enumerate}

\begin{enumerate}
\def\labelenumi{\alph{enumi})}
\item
  Show that the set V of \(x \in E\) such that \(\Phi(x, x)\) is of the form \(\alpha + \bar{\alpha}\) is a vector subspace of E.
\item
  Let \(V_{1}=V\cap V^{0}\) , \(q=\dim V_{1}\) , \(V_{2}\) a supplement of \(V_{1}\) with respect to V, \(V_{3}\) a supplement of \(V_{1}\) with respect to \(V^{0}\) . Show that there exists a basis \((e_{i})_{1\leqslant i\leqslant2q}\) of \((\mathrm{V}_{2}+\mathrm{V}_{3})^{0}=\mathrm{V}_{2}^{0}\cap\mathrm{V}_{3}^{0}\) such that the vectors \(e_1, \ldots, e_q\) form a basis of \(V_1\) and that one has \(\Phi(e_i, e_{q+j}) = \delta_{ij}\) for \(1 \leqslant i \leqslant q\) , \(1 \leqslant j \leqslant q\) .
\item
  Let \(\mathbf{G}(\Phi)\) be the group of metric automorphisms of E (for \(\Phi\) ). Show that for every \(u \in \mathbf{G}(\Phi)\) , one has \(u(x) = x\) for all \(x \in V^0\) .
\item
  For every \(u \in \mathbf{G}(\Phi)\) , one has \(u(V) = V\) ; let \(u_{v}\) be the restriction of \(u\) to \(V\) , and let \(G_{v}\) be the group formed by the \(u_{v}\) . Show that: 1° the kernel of the homomorphism \(u \to u_{v}\) of \(\mathbf{G}(\Phi)\) onto \(G_{v}\) is commutative; 2° if \(\Phi_{2}\) is the restriction of \(\Phi\) to \(V_{2}\) and \(\mathbf{G}(\Phi_{2})\) the group of metric automorphisms of \(V_{2}\) for \(\Phi_{2}\) , there exists a homomorphism from \(G_{v}\) onto \(\mathbf{G}(\Phi_{2})\) whose kernel is commutative (use \(b\) ) and \(c\) ).
\item
  We assume that A is commutative and J is the identity; let E be a vector space of dimension 3 over A, \((e_{i})_{1\leqslant i\leqslant3}\) a basis of E, \(\Phi\) the nondegenerate symmetric form on E whose matrix with respect to \((e_{i})\) is
\end{enumerate}

\[
\left( \begin{array}{c c c} 1 & 0 & 0 \\ 0 & 0 & 1 \\ 0 & 1 & 0 \end{array} \right).
\]

Show that all isotropic vectors in E are contained in the hyperplane spanned by \(e_{2}\) and \(e_{3}\) (cf. exercise 4(a)). Give an example of a non-isotropic plane containing only one isotropic line (cf. exercise 2(a)). Show that there exists no automorphism \(u \in \mathbf{G}(\Phi)\) such that \(u(e_{1}) = e_{1} + e_{2}\) , even though one has \(\Phi(e_{1}, e_{1}) = \Phi(e_{1} + e_{2}, e_{1} + e_{2})\) .

